\documentclass[a5paper,10pt,twoside,openright]{memoir}
\usepackage{multicol, multirow, array}
\usepackage{fontspec}
\usepackage{anyfontsize}
\usepackage[pagecolor=none,dvipsnames]{xcolor}
\usepackage{ragged2e}
\usepackage{amsmath}
\usepackage{amssymb}
\usepackage[hidelinks]{hyperref}
\usepackage{url}
\usepackage[margin=0.8in]{geometry}
\usepackage{float, hhline}
\usepackage{graphicx}
\usepackage{booktabs}
\usepackage{textcomp}
\usepackage{expex}
\usepackage{enumitem}
\usepackage[calc,english]{datetime2}
\usepackage{suffix}

%-----CONFIGURATION------
%------------------------

\setmainfont{Charis SIL}[CharacterVariant=43:1]
\restylefloat{table}

\setsecnumdepth{subsubsection}
\settocdepth{subsubsection}

\lingset{glstyle=nlevel,numoffset=3em,textoffset=1.5em,exskip=.75ex,belowglpreambleskip=.25ex,aboveglftskip=.25ex}

\DTMnewdatestyle{eurodate}{%
    \renewcommand{\DTMdisplaydate}[4]{%
        \number##3.\nobreakspace%           day
        \DTMmonthname{##2}\nobreakspace%    month
        \number##1%                         year
    }%
    \renewcommand{\DTMDisplaydate}{\DTMdisplaydate}%
}

\DTMsetdatestyle{eurodate}

\renewcommand{\arraystretch}{1.4}

\lingset{belowpreambleskip=2ex}

%-------COMMANDS---------
%------------------------

\newcommand{\lang}{ɁaɁa-\textsc{f}a\textsc{m}}
\newcommand{\longv}{ː}
\newcommand{\sqbrack}[1]{$\langle$#1$\rangle$}
\newcommand{\phipa}[1]{/#1/}
\newcommand{\bripa}[1]{[#1]}
\newcommand{\ttilde}{\raise.17ex\hbox{$\scriptstyle\sim$}}
\newcommand{\rootpart}{$\Theta$}
\newcommand{\glotstop}{ʔ}
\newcommand{\bigglot}{Ɂ}
\newcommand{\lilglot}{ɂ}
\newcommand{\nm}{\symbol{"2205}}
\newcommand{\tiebar}{͡}
\newcommand{\famwordold}[5]{#1\textsc{#2}#3\textsc{#4}#5}
\newcommand\famword[1]{{\addfontfeatures{Letters=UppercaseSmallCaps}#1}}
\newcommand{\famq}[1]{»#1«}
\newcommand{\glem}[1]{\underline{\smash{#1}}}

\newenvironment{conversation}
    {\begin{table}[ht]
        \begin{tabular}{p{0.962\textwidth}}
        \toprule
        \textbf{Conversation} \\
        \midrule
    }
    { 
        \bottomrule
        \end{tabular}
    \end{table}
    }

\newcommand{\diagline}[2]{\hspace{2em} \textbf{\uppercase{#1:}} {#2} \\}

%-------TITLE PAGE-------
%------------------------

\title{\fontsize{70}{70}\selectfont \itshape \famword{iFaaMak!} \\ \fontsize{25}{30} \sffamily A Beginner's Guide to Narish}
\author{Bethany E. Toma, Knut F. K. Ulstrup}
\date{}%\today}

%--------MAIN DOC--------
%------------------------

\begin{document}

\begin{titlingpage}
    \maketitle
    \begin{figure}[ht]
        \includegraphics[width=\textwidth]{nareland3small.png}
    \end{figure}
\end{titlingpage}

\frontmatter

\begin{quote}
    \textbf{\famq{\famword{\Large NaJDa FaM baj iFaaM se BaBauru \\ JaLa FaM baj iFaaM se TaLWuru \\ da TaPa FaM baj se iFaaM se mi FaRaariibi.}}}

    \textit{\small ``Speaking a single language is easy, \\ speaking many languages is a burden, \\ but speaking all languages sets you free.'' \\ -Old Narish proverb}
\end{quote}

\newpage

\tableofcontents

\mainmatter

\chapter{Introduction to Narish}

\chapter{The Narish alphabet}

Narish uses the latin alphabet, slightly modified to better suit the language.

\begin{table}[ht]
    \centering
    \begin{tabular}{ccccl}
        \toprule
        \multicolumn{3}{l}{\footnotesize letter cases} & \footnotesize sound & \footnotesize example \\
        \midrule
        \lilglot & \bigglot & \bigglot & /ʔ/ & like the break in `uh-oh' \\
        a & A & \textsc{a} & /a/            & h\glem{a}rd \\
        b & B & \textsc{b} & /p{\ttilde}b/  & \glem{b}ottle \\
        c & C & \textsc{c} & /c{\ttilde}c{\tiebar}ç/ & \glem{ch}ance \\
        d & D & \textsc{d} & /t{\ttilde}d/  & \glem{d}onkey \\
        e & E & \textsc{e} & /ə/            & salm\glem{o}n \\
        f & F & \textsc{f} & /f/            & \glem{f}acet \\
        h & H & \textsc{h} & /h{\ttilde}χ/  & \glem{h}appy, lo\glem{ch} \\
        i & I & \textsc{i} & /i/            & f\glem{ee}l \\
        j & J & \textsc{j} & /j/            & \glem{y}awn \\
        k & K & \textsc{k} & /k{\ttilde}k{\tiebar}x/ & \glem{c}ar \\
        l & L & \textsc{l} & /l/            & be\glem{l}ieve \\
        m & M & \textsc{m} & /m/            & \glem{m}an \\
        n & N & \textsc{n} & /n/            & \glem{n}otion \\
        p & P & \textsc{p} & /pʰ{\ttilde}p{\tiebar}ɸ/& \glem{p}allet \\
        q & Q & \textsc{q} & /q/            & like /k/ further back \\ 
        r & R & \textsc{r} & /ɾ{\ttilde}ɹ/  & bu\glem{tt}er, \glem{r}eally \\
        s & S & \textsc{s} & /s/            & \glem{s}ing \\
        t & S & \textsc{t} & /tʰ{\ttilde}t{\tiebar}s/& \glem{t}ime \\
        u & U & \textsc{u} & /u/            & h\glem{u}l\glem{u} \\
        w & W & \textsc{w} & /w/            & \glem{w}ater \\
        y & Y & \textsc{y} & /j/            & \glem{y}awn \\
        \bottomrule
    \end{tabular}
\end{table}

\chapter{Forming words}

The peculiar writing style will seem odd to most, but it is this way for good reason. In many words, certain letters are more important than others as they point to a common overarching meaning, called the \emph{root}, and it's important that these letters stand out for readability.

In the words \famword{KarLi} \textit{`water'}, \famword{KaLa} \textit{`wet'}, \famword{iKaaL} \textit{`it's raining'}, and \famword{aKiLu} \textit{`bottle'}, the letters \textsc{k} and \textsc{l} are part of the root \emph{K-L}, which doesn't have a meaning of its own, but clearly ties together all these words related to water in some way. All of these words are comprised of two root consonants and a template called a \emph{stem} that these consonants slot into. 

\chapter{Saying hello}

Of course, the first part of any conversation is greetings. It's important that you convey greetings and wishes properly, as their exact forms are determined by your gender expression.

\begin{conversation}
    \diagline{A}{\famword{FaSajara}.}
    \diagline{B}{\famword{FaSunara}.}
\end{conversation}

In the above example, person A identifies themselves as male by using the male form \emph{\famword{FaS\glem{aj}ara}}, and person B identifies themselves as female by means of the female form \emph{\famword{FaS\glem{un}ara}}. When you say hello, you're also letting the other person know exactly how you identify, and they're also telling you back! There are in total four different forms, as listed below:

\begin{table}[ht]
    \centering
    \begin{tabular}{rc}
        \toprule
        Male & \famword{FaSajara} \\
        Female & \famword{FaSunara} \\
        Nonbinary & \famword{FaSujara} \\
        Unspecified & \famword{FaSanara} \\
        \bottomrule
    \end{tabular}
\end{table}

\end{document}